% SOURCE_AUTHORITY: sga5_fr_workpass.tex, lines 12870--12894 (LF-normalized unit).
% SOURCE_SHA256: 791F4EFFC5E02832D5D77ED03518C8156D6F07E4C8238B03545DB93D883FBB28
De fato, o objeto de $\Delta$ definido por $\Hom^\bullet_{\Lambda[G]}(Sw_x^\Lambda,M)$ não depende, salvo isomorfismo, do recobrimento de Galois $C'\to C$ que trivializa $F|U$. Para demonstrá-lo, basta comparar os elementos associados a dois recobrimentos
\[
C'\xrightarrow{p}C,\qquad C''\xrightarrow{p'}C'\xrightarrow{p}C.
\]
Neste caso, se $G'$ é o grupo de Galois do recobrimento $C''\xrightarrow{p''}C$, então $G$ é um grupo quociente de $G'$. Seja
\[
\varphi:\mathbb Z_\ell[G']\longrightarrow \mathbb Z_\ell[G]
\]
o epimorfismo canônico associado. Os módulos de Swan estão então ligados pela relação
\[
Sw_{G,x}\simeq Sw_{G',x}\otimes_{\mathbb Z_\ell[G']}\mathbb Z_\ell[G],
\]
que resulta imediatamente do fato de que o caráter $a_{G,x}$ é o caráter induzido de $G'$ a $G$ ([3] VI Prop. 3).

Por outro lado, tem-se um isomorfismo de complexos de $A[G']$-módulos
\[
F_{\bar\eta''}\simeq [F_{\bar\eta'}]_{(\varphi)}.
\]
Disso resulta um isomorfismo canônico de complexos de $A$-módulos entre
\[
\Hom^\bullet_{\Lambda[G]}(Sw_{G,x}^\Lambda,F_{\bar\eta})
\quad\text{e}\quad
\Hom^\bullet_{\Lambda[G']}(Sw_{G',x}^\Lambda,F_{\bar\eta''})
\]
(fórmula de adjunção).
